\section{Benchmark Design}\label{sec:benchmark}

\subsection{Query Taxonomy}\label{sec:query-taxonomy}

We define five query types (T1--T5), each targeting a different aspect of
knowledge retrieval capability. Table~\ref{tab:query-types} summarizes the
taxonomy.

\begin{table}[htbp]
\centering
\caption{Query type taxonomy with examples and ground truth sources.}
\label{tab:query-types}
\begin{tabular}{lp{3.5cm}p{4.5cm}l}
\toprule
\textbf{Type} & \textbf{Description} & \textbf{Example} & \textbf{Ground Truth} \\
\midrule
T1 & Entity lookup & ``What is Composite Function?'' & ConceptLabel + TaxonomyID \\
T2 & Direct dependency & ``What are prerequisites for Composite Function?'' & Dependencies column \\
T3 & Multi-hop path & ``What is the chain from Function to Taylor Series?'' & BFS path in DAG \\
T4 & Category aggregate & ``List all FOUND concepts'' & Filter by TaxonomyID \\
T5 & Cross-concept & ``How does Domain and Range relate to Inverse Function?'' & Shared neighbors \\
\bottomrule
\end{tabular}

\vspace{4pt}
\noindent\small\textbf{Note on T1 (entity lookup):} T1 queries ask for a concept explanation
(``What is X?''). CKG contains no explanatory prose---it can return only the
concept's TaxonomyID and dependencies in response. T1 is therefore a
\emph{RAG-favorable} query type deliberately included to test the boundary of
CKG's capability and to confirm that the benchmark is not constructed to favor
CKG universally. T2--T4 queries have structural ground truth derived from DAG
edges; T1 does not.
\end{table}

\subsection{Query Generation}\label{sec:query-gen}

Queries are auto-generated from each domain's CSV using
\texttt{generate\_queries.py} with a fixed random seed of 42.

Per domain: $\sim$175 queries (50 T1 + 50 T2 + 25 T3 + 12 T4 + 38 T5).
Total: $\sim$4,375 queries across 25 domains.

\begin{itemize}
  \item \textbf{T1}: Random sample of 50 concepts; query is ``What is \{label\}?''
  \item \textbf{T2}: Random sample of 50 concepts with $\geq$1 dependency
  \item \textbf{T3}: Random pairs of foundational/terminal concepts with path length 2--5
  \item \textbf{T4}: One query per taxonomy category
  \item \textbf{T5}: Random sample of 38 directly connected concept pairs
\end{itemize}

\subsection{Ground Truth Validity}\label{sec:validation}

Benchmark ground truth is derived \emph{deterministically} from DAG edges:
T2 answers are the direct dependency labels of a concept; T3 answers are the
BFS shortest-path node sequence between two concepts; T4 answers are all
concept labels sharing a TaxonomyID; T5 answers are the union of BFS path
nodes between a randomly selected concept pair. Because derivation is
algorithmic, inter-annotator $\kappa$ does not apply to the ground truth
generation process---there are no annotators to disagree.

Benchmark validity is instead assessed structurally. The T1 negative-control
result (\textbf{CKG F1 $= 0.207$} on entity lookup) confirms the evaluation
is not constructed to favor CKG universally: a system reading only DAG edges
appropriately fails on definitional queries the DAG cannot answer. The
parallel RAG and GraphRAG T1 scores (0.094 and 0.108 respectively) show that
T1 is hard for all systems, ruling out an inflated CKG baseline. Additionally,
the low GraphRAG T4 score (0.054 vs.\ CKG 0.964) confirms that the taxonomy
advantage is structural---not an artifact of prompting---since GraphRAG and
CKG use identical prompts and the same LLM.

\subsection{Reproducibility Protocol}\label{sec:reproducibility}

\begin{itemize}
  \item All systems use Claude Sonnet 4.6 at temperature$=$0.
  \item Token counts via Anthropic \texttt{count\_tokens()} API.
  \item 3 runs per query, variance reported.
  \item Fixed random seed: 42.
  \item Benchmark version locked: v1.0.0.
  \item One-command reproduction: \texttt{python evaluation/harness.py --reproduce-table-1}
\end{itemize}
